% Transcription — fonds Alexandre Grothendieck, shelfmark 80, batch 1 (pages 1-20).
% Established from the facsimile archives/batches/80/batch-01.pdf, cut from a
% copy of 80.pdf fetched through the project's own relay
% (https://grothendieck-relay.onrender.com/source/80.pdf), not uploaded by the
% user; PDF page k = archivists' page k, no cover sheet. Per the skill
% transcribe-grothendieck. The folder is seventy-two pages in four batches
% (1-20, 21-40, 41-60, 61-72), transcribed in parallel by four passes; this is
% the first, and no neighbouring file was relied on.
%
% Pass: Opus 5.5 (claude-opus-5-5), 2026-09-26 — first pass, unchecked against the pages by a human.
%
% Dating and title. \dating is the folder's own date, copied verbatim from
% src/content/catalogue.ts; the folder belongs to the inventory group
% « Géométrie et topologie combinatoire » (69-89), but since it has a date of
% its own the group's range is not added. Title from the catalogue, nothing
% bracketed (trailing full stop dropped, as in folders 75 and 86).
%
% Pages skipped: 1, 8, 15 and 19, blank cream separator sheets with nothing
% on them. Pages summarised: none. Pages given only as a figure described in
% a French \note, with their few formulas and labels set: 9, 11, 12, 13, 14
% (and the drawings of 20). Pages transcribed from his hand: 2-7 (ink on
% computer-listing paper, the printout showing through), 10, 16-18, 20 (ink on
% white paper).
%
% His pagination: pages 2-7 form one numbered run, 1 (circled, p. 2), 2
% (circled, p. 3), 3) (p. 4), 4) (p. 5), 5) (p. 6), 6) (p. 7); recorded in a
% \note at each page. Nothing else in the batch carries his numbers. Page 20
% numbers its drawings (1)-(4) and a fifth configuration (5).
%
% What the batch is about.
%   2-7   « Classifications pseudo-droites » (pencil heading on p. 2). For
%         I of card n+1, i_0 in I, I_0 = I \ {i_0}: sets C(I) ≃ C_aff(I_0),
%         C(n+1) = C(I)/S_I, C_aff(n) = C_aff(I_0)/S_{I_0}, C(n) of
%         iso classes of pseudo-line systems (X, Σ, Δ), with forgetful maps
%         α, β, β' and their fibres (Σ/G_{X,Σ}; for n = 5 fibres of card 1
%         and 2); the « stable » variants C_st, C_aff st. P. 3: a stable class
%         ξ gives I_0 a polygonal structure and a double cover Ĩ_0(ξ); in
%         terms of the disc D = Déc(X, Δ) with a stable system of chords
%         Σ̃ indexed by I_0, Ĩ_0(ξ) is the set of endpoints on ∂D. For a fixed
%         combinatorial polygon S of order 2n, C̃_st(S) (S-pinned stable chord
%         systems) maps to C^!_aff st(I_0), I_0 = S/antipody, with fibres of
%         card 2 (pp. 3-4). Pp. 5-6: counting, Card α^{-1} = (n!/2n) Card
%         α^{!-1}, card α^{!-1}(ξ) = 2n/N(ξ), Card α̃^{-1} = 4n/N(ξ) with
%         N(ξ) = card Aut(X, Σ, Δ); examples n = 3 (Aut ≃ S_3), n = 4
%         (card C̃_st(S) = 8), n = 5 (6 classes); p. 7 « Résumé ».
%   9     A large drawing of an arrangement in a disc (Δ_A..Δ_C', D_A..D_C,
%         regions labelled by number of sides) and a count of triangles (7),
%         quadrilaterals (24), pentagons (6).
%   10    Unrelated to pseudo-lines: Stein factorisation — for Z' -> X'
%         steinian, a unique X' -> Z̃; the initial object of the category of
%         steinian arrows above Z -> X; compatibility with base change,
%         conditions a) Z -> X compact ?, b) X' -> X loc. O-acyclic.
%   11-14 Drawings: polyhedra (octahedron, cube, truncated octahedron),
%         48 = 2·4!, 12 = 8·3·2/4; a six-line arrangement with a central
%         pentagon; a construction with inequalities a+a' ≤ 2, b+b' ≥ 2,
%         a'+b > 2; an arrangement D_0..D_4 with shaded half-planes.
%   16-18 « Système d'hyperplans (projectifs) réels » (his heading): for
%         hyperplanes H_i of a real V, D_H = {H, V'_H, V''_H} ≃ F_3,
%         φ : S = V*/R*_+ -> E = ∏ D_{H_i}, Σ = φ(S); closures
%         Z̄*_α = ∪_{β ≤ α} Z*_β for the order 0 < α', α'' on each factor;
%         Σ indexes a stratification of the sphere satisfying the frontier
%         axiom, with convex strata; a cancelled paragraph on Inf; Σ open in E.
%   20    Arrangements of 4, 5, 5, 6 pseudo-lines, the claim that every
%         arrangement not reduced to 4 lines contains a configuration of type
%         (2), (3) or (4), and an NB on type (5), 7 lines (quadrilateral and
%         diagonals, reading doubtful). Ends on a struck « Si elle ne contient ».
%
% Notation fixed once for the batch: the subscript « aff » as he writes it (a
% small a with doubled ascenders, read as aff), « st » for stable; \mathfrak{S}
% for his gothic S (symmetric group); \mathbb{D} for his open-faced D (dihedral
% group on pp. 5-7, the three-element D_H on p. 17); \xi^{q} for the small
% superscript on ξ in pp. 5-7 (reading doubtful, said in a note); \mathbf{S}
% for the doubled S of p. 16, \mathbf{E}, \mathbf{F}_3.
%
% Readings to check first: the subscript « aff »; the superscript of ξ^q;
% the arrow labels of the diagram on p. 2; the marginal notes on pp. 16-17
% (mostly \ill{}); the last lines of p. 18 and the NB of p. 20.

\documentclass[11pt,a4paper]{article}
% Le préambule commun à toutes les transcriptions du fonds Grothendieck.
%
% Il tient en une idée : **l'incertitude se compose, elle ne se résout pas.**
% Une transcription qui lisse un mot douteux a détruit la seule chose pour
% laquelle on la faisait — un lecteur doit pouvoir savoir, à chaque endroit,
% ce qui était sur le feuillet et ce qui a été ajouté. Les six macros qui
% suivent sont donc l'appareil critique, et non de la décoration.
%
% Les mêmes commandes sont reconnues par `scripts/render.mjs`, qui produit la
% vue de lecture du site. Une macro ajoutée ici doit l'être aussi là-bas,
% faute de quoi le rendu échouera bruyamment — ce qui est voulu : un rendu
% silencieusement faux est pire qu'un rendu absent.

\ProvidesPackage{grothendieck}[2026/08/08 Transcriptions du fonds Grothendieck]

\usepackage{amsmath,amssymb,amsthm}
\usepackage{mathrsfs}
\usepackage{stmaryrd}
% Les diagrammes commutatifs sont partout dans ce fonds : c'est la langue
% même de la géométrie algébrique de Grothendieck, et une transcription qui
% les rendrait en prose ne transcrirait rien.
\usepackage{tikz-cd}
% Français par défaut — c'est la langue des feuillets — mais l'anglais est
% chargé aussi : la traduction et le résumé sont des documents anglais, et la
% typographie française (espaces avant les deux-points) y serait une faute.
% Ces documents basculent par \selectlanguage{english} en tête de corps.
\usepackage[english,french]{babel}
\usepackage{fontspec}
\usepackage{geometry}
\geometry{a4paper,margin=25mm,marginparwidth=28mm}
\usepackage{marginnote}
\usepackage{xcolor}
% ulem, not soul. `soul`'s \st and \ul are built on a character-by-character
% scan that XeLaTeX's font machinery breaks: the document fails with an
% undefined control sequence rather than degrading. ulem does the same two jobs
% and is Unicode-engine safe. `normalem` stops it hijacking \emph, which the
% transcriptions use for Grothendieck's own emphasis.
\usepackage[normalem]{ulem}
\usepackage{hyperref}
\hypersetup{colorlinks=true,linkcolor=black,urlcolor=black,pdfborder={0 0 0}}

% --- Métadonnées du lot -----------------------------------------------------
% Reprises telles quelles par le script de rendu, qui les affiche en tête de
% la vue de lecture. Elles fixent ce que le fichier prétend couvrir : sans
% elles, une transcription flotte sans point d'attache dans un fonds de seize
% mille feuillets.

\newcommand{\folder}[1]{\gdef\grothFolder{#1}}
\newcommand{\batch}[1]{\gdef\grothBatch{#1}}
\newcommand{\leaves}[2]{\gdef\grothFirst{#1}\gdef\grothLast{#2}}
\newcommand{\foldertitle}[1]{\gdef\grothFoldertitle{#1}}
\gdef\grothFolder{?}\gdef\grothBatch{?}\gdef\grothFirst{?}\gdef\grothLast{?}\gdef\grothFoldertitle{}

% --- Le résumé --------------------------------------------------------------
%
% Il ouvre la lecture modernisée, sous le titre « Résumé », et il est composé
% en petit. Ce n'est pas un ornement : c'est le seul passage du document qui
% ne soit pas des mathématiques, et un lecteur doit pouvoir le reconnaître
% avant de l'avoir lu — puis le sauter s'il connaît déjà le sujet, ou s'y
% arrêter s'il ne le connaît pas. Le filet qui le suit marque la reprise du
% corps.

\newenvironment{resume}
  {\small\setlength{\parskip}{0.45em}}
  {\par\medskip\hrule\medskip}

% --- Mots-clés ---------------------------------------------------------------
%
% En anglais, en clôture du résumé de la lecture modernisée : le vocabulaire
% moderne sous lequel chercher ce que les feuillets construisent. Le manifeste
% du site (`npm run manifest`) les extrait pour étiqueter la cote entière —
% cette ligne est la source unique des tags, il n'y a pas de fichier à côté.
\newcommand{\keywords}[1]{\par\smallskip\noindent
  {\footnotesize\textsc{Keywords} --- \textit{#1}}\par}

% --- L'appareil critique ----------------------------------------------------

% Le numéro de feuillet, en marge. C'est l'ancre qui permet de régler un
% désaccord sur une formule en regardant *un* feuillet plutôt qu'une chemise
% de six cents. Le site s'en sert aussi pour tourner les pages du fac-similé
% quand on fait défiler la transcription.
\newcommand{\leaf}[1]{%
  \marginnote{\footnotesize\color{gray!70}\textsf{#1}}%
  \label{leaf:#1}\ignorespaces}

% Illisible : on ne devine pas. Un mot inventé qui se lit comme les autres est
% le pire résultat possible.
\newcommand{\ill}{\textcolor{red!60!black}{[\,\ldots\,]}}

% Lecture douteuse : le mot est donné, et signalé comme incertain.
\newcommand{\uncertain}[1]{\uline{#1}}

% Ajout de l'éditeur — une lettre manquante, un mot sous-entendu.
\newcommand{\add}[1]{\textcolor{blue!50!black}{[#1]}}

% Biffé par Grothendieck lui-même. Ce qu'il a rayé fait partie du document :
% une rature dit souvent où la pensée a changé de direction.
\newcommand{\struck}[1]{\sout{#1}}

% Note du transcripteur, sur l'état matériel du feuillet.
\newcommand{\note}[1]{\par\smallskip{\small\color{gray!60!black}\textit{[#1]}}\par\smallskip}

% Note marginale de Grothendieck, distincte du corps du texte.
\newcommand{\marginal}[1]{\par\smallskip\noindent\hspace{1em}%
  {\small\color{gray!60!black}#1}\par\smallskip}

% --- Titre ------------------------------------------------------------------

\AtBeginDocument{%
  \begin{center}
    {\large\bfseries Fonds Alexandre Grothendieck --- cote n\textsuperscript{o}~\grothFolder}\\[2pt]
    {\itshape\grothFoldertitle}\\[4pt]
    {\small feuillets \grothFirst--\grothLast\ (lot \grothBatch)}\\[2pt]
    {\footnotesize\color{gray!60!black}Transcription établie d'après le fac-similé numérisé
     par l'Université de Montpellier.\\
     Les lectures incertaines sont soulignées, les illisibles notées [\,\ldots\,],
     les ajouts entre crochets.}
  \end{center}
  \vspace{1em}}

\endinput


\folder{80}
\batch{1}
\pages{1}{20}
\dating{[à partir de 1981]}
\watermark{Édition de démonstration}
\foldertitle{Systèmes de pseudo-droites : notes manuscrites (s.d.)}

\begin{document}

\section{Classifications pseudo-droites}
\note{titre écrit au crayon, de sa main, en tête de la page 2, au-dessous du premier diagramme}

\page{2}
\note{numéroté \emph{1} (entouré) en haut à droite, de sa main ; la suite numérotée de 1 à 6 court sur les pages 2 à 7}
$I$, $i_0 \in I$, $\operatorname{card} I = n+1$, $I_0 = I \setminus \lbrace i_0 \rbrace$.

\[
\mathrm{C}(I) \simeq \mathrm{C}_{\mathrm{aff}}(I_0) \longrightarrow \mathrm{C}_{n+1} \simeq \mathrm{C}(I)/\mathfrak{S}_I ,
\qquad
\mathrm{C}_{\mathrm{aff}}(n) \simeq \mathrm{C}_{\mathrm{aff}}(I_0)/\mathfrak{S}_{I_0} \simeq \mathrm{C}_{\mathrm{aff}}(I)/\mathfrak{S}_I
\]
\note{ce premier diagramme (une flèche horizontale et une flèche verticale de $\mathrm{C}(I)$ vers $\mathrm{C}_{\mathrm{aff}}(n)$), avec les données $I$, $i_0$, $I_0$ à droite, est biffé par de longs traits obliques ; il est repris plus bas. L'indice « $\mathrm{aff}$ » est \uncertain{lu} ainsi : il l'écrit comme un petit « a » suivi de deux hampes doublées, et on le garde tel pour tout le lot. $\mathfrak{S}$ rend son S gothique (groupe symétrique)}

\begin{tikzcd}[column sep=small, row sep=large, nodes={font=\scriptsize}]
\mathrm{C}(I) \simeq \mathrm{C}_{\mathrm{aff}}(I_0) \arrow[dr, "\alpha"'] & & \\
& \mathrm{C}_{\mathrm{aff}}(n) \simeq \mathrm{C}_{\mathrm{aff}}(I_0)/\mathfrak{S}_{I_0} \simeq \mathrm{C}_{\cdots}(I)/\mathfrak{S}_{I_0} \arrow[dl, "\beta"'] \arrow[dr, "\beta'"] & \\
\mathrm{C}(n+1) \simeq \mathrm{C}(I)/\mathfrak{S}_I & & \mathrm{C}(n)
\end{tikzcd}

\note{étiquettes des flèches : $\alpha$ = « oubli du $I$-épinglage sauf $\Delta_{i_0} = \Delta$ », la flèche barrée de deux petits traits ; $\beta$ = « \struck{\ill{}} oubli du $\Delta$ \ill{} \ill{} $\Delta$ », avec un signe $\simeq$ couché le long de la flèche ; dans $\mathrm{C}_{\cdots}(I)$ l'indice est biffé et illisible ; $\beta'$ = « \struck{oubli de $\Delta$} gommage de $\Delta$ ». À droite du diagramme : « classes d'iso de systèmes $(X, \Sigma, \Delta)$ : $(X, \Sigma)$ syst.\ de ps dr.\ et $\Delta \in \Sigma$ »}

$\beta$ surjection, la fibre en \struck{$X$}, $\mathrm{cl}_{n+1}(X, \Sigma)$ est l'ens.\ des types d'éléments de $\Sigma$, i.e.\ $\Sigma / G_{X,\Sigma}$

[Si $n = 5$, $n+1 = 6$, il y a deux fibres qui sont de card $1$, deux de card.\ $2$

$\beta'$ surjective (?)

La fibre de $\mathrm{cl}_n(X, \Sigma^{*})$ est l'ens des types de positions relatives p.r.\ : $(X, \Sigma^{*})$

[Si $n = 5$, $\mathrm{C}(n)$ \uncertain{a un} élément, la fibre est de card.\ $6$]
\note{les deux crochets sont écrits l'un en face de l'autre, dans les deux colonnes ; le premier n'est pas refermé}

\emph{NB} $\mathrm{C}_{\mathrm{aff}\,\mathrm{st}}(I_0) \subset \mathrm{C}_{\mathrm{aff}}(I_0)$, avec $\mathrm{C}_{\mathrm{aff}\,\mathrm{st}}(I_0) = \mathrm{C}_{\mathrm{st}}(I)$ et $\mathrm{C}_{\mathrm{aff}}(I_0) = \mathrm{C}(I)$,
stable par $\mathfrak{S}_I$, a fortiori par $\mathfrak{S}_{I_0}$, prenant \uncertain{les} quotients \ill{} trouve \ill{} les diagrammes analogues \ill{}
$\mathrm{C}_{\mathrm{st}}(I) = \mathrm{C}_{\mathrm{aff}\,\mathrm{st}}(I_0)$, $\mathrm{C}_{\mathrm{aff}\,\mathrm{st}}(n)$, $\mathrm{C}_{\mathrm{st}}(n)$, $\mathrm{C}_{\mathrm{st}}(n+1)$.
\note{les égalités $\mathrm{C}_{\mathrm{aff}\,\mathrm{st}}(I_0) = \mathrm{C}_{\mathrm{st}}(I)$ et $\mathrm{C}_{\mathrm{aff}}(I_0) = \mathrm{C}(I)$ sont écrites verticalement, par des « $=$ » sous chaque terme de l'inclusion}

\page{3}
\note{numéroté \emph{2} (entouré), en haut à gauche}
Considérons un élément de
\[
\mathrm{C}_{\mathrm{st}}(I) \simeq \mathrm{C}_{\mathrm{aff}\,\mathrm{st}}(I_0)
\]
$\xi = \mathrm{cl}(X, \Sigma = (\Delta_i)_{i \in I})$. Alors $I_0$ hérite d'une \emph{structure polygonale} (dépendant de $\xi$) (\ill{} \ill{} \uncertain{card} $I_0 \geqslant 3$)
\note{« (dépendant de $\xi$) » est ajouté au-dessus de la ligne et renvoyé par un trait sous « structure polygonale »}

De plus, on trouve un polygône $\widetilde{I}_0(\xi)$ qui est un revêtement à deux feuillets de $I_0(\xi)$.
\[
\widetilde{I}_0(\xi) \longrightarrow I_0(\xi)
\]
\note{la flèche est verticale dans la marge gauche, en face des lignes qui suivent}
Si on \struck{explicite} traduit la structure $(X, \Sigma = (\Delta_i)_{i \in I})$, en mettant : par $\Delta = \Delta_{i_0}$ : à l'aide d'un disque $\mathbf{D} = \mathrm{Déc}(X, \Delta)$, \struck{alors} et d'un système \add{stable} de cordes $\widetilde{\Sigma}$ dans $\mathbf{D} = \widetilde{X}$ indexé par $I_0$, on trouve \struck{\ill{} avec des sommets sur $\partial \mathbf{D}$.}
$\widetilde{I}_0(\xi) \simeq{}$ ens des sommets sur $\partial\mathbf{D} = \partial\widetilde{X}$
(rem.\ : deux feuillets de $\Delta$)
\note{« stable » est ajouté au-dessus de « système de cordes », souligné ; « traduit » est écrit au-dessus de « explicite » biffé. Son $\Delta$ porte un trait intérieur, comme s'il s'agissait d'un $\Delta$ ouvert d'un côté ; on le rend par $\Delta$ partout. « $\mathrm{Déc}$ » est son abréviation, gardée telle}

Soit maintenant $S$ un ens.\ polygonal \uncertain{de} \uncertain{même} ordre $2n$, donné une fois pour toutes. On s'intéresse aux structures de syst.\ stables de cordes $S$-épinglées (mod $S$-isom.), soit $\widetilde{\mathrm{C}}_{\mathrm{st}}(S)$ l'ens des cl.\ de ces structures. Posons
$I_0 = S / a_S$, \quad $a_S$ antipodisme de $S$.
On a une application naturelle
\[
\widetilde{\mathrm{C}}_{\mathrm{st}}(S) \longrightarrow \mathrm{C}_{\mathrm{aff}\,\mathrm{st}}(I_0, \ldots)
\]
\note{le second argument, après la virgule, est raturé et illisible}
\struck{dont la fibre en $\xi$ est un cov.\ \ill{} \ill{} l'ens \ill{}}

\page{4}
\note{numéroté \emph{3)} en haut à gauche : la page continue la phrase interrompue au bas de la page 3}
dont l'image est le sous-ens
\[
\mathrm{C}^{!}_{\mathrm{aff}\,\mathrm{st}}(I_0) \overset{\mathrm{def}}{=} \mathrm{C}_{\mathrm{aff}\,\mathrm{st}}(I_0, \pi_0) \subset \mathrm{C}_{\mathrm{aff}\,\mathrm{st}}(I_0)
\]
formé des $\xi$ tels que la structure polygonale correspondante sur $I_0$ soit la structure polygonale $\pi_0$ obtenue sur $I_0$, déduite de la structure polygonale de $S$ par passage au quotient. Donc on a
\[
\widetilde{\mathrm{C}}_{\mathrm{st}}(S) \longrightarrow \mathrm{C}^{!}_{\mathrm{aff}\,\mathrm{st}}(I_0)
\]
où le ${}^{!}$ indique que $I_0$ est muni de $\pi_0$, i.e.\ qu'on prend $\mathrm{C}_{\mathrm{aff}\,\mathrm{st}}(I_0, \pi_0)$. La \struck{fibre de} fibre de $\xi \in \mathrm{C}^{!}_{\mathrm{aff}\,\mathrm{st}}(I_0)$ s'identifie à l'ens des iso polygonaux
\[
S \xrightarrow{\ \sim\ } \widetilde{I}_0(\xi)
\]
qui \struck{sont} rendent commutatif le triangle

\begin{tikzcd}
S \arrow[rr, "\sim"] \arrow[dr, "\mathrm{can}"'] & & \widetilde{I}_0(\xi) \arrow[dl, "\mathrm{can}"] \\
& I_0 &
\end{tikzcd}

il est de cardinal $2$ ; \emph{on passe de l'un à l'autre par antipodisme}
\note{au début de la page, un renvoi en tête insère $\mathrm{C}^{!}_{\mathrm{aff}\,\mathrm{st}}(I_0)$ au-dessus de la ligne, et un sigle raturé précède la formule ; nous lisons la notation $\mathrm{C}^{!}$ d'après la suite de la page, où le « ! » est expliqué}

\struck{\emph{Ex} Card $I_0 = 5$, $\mathrm{C}_{\mathrm{aff}\,\mathrm{st}}(I_0) \simeq \mathrm{C}_{\mathrm{st}}(I_0 \sqcup \lbrace i_0 \rbrace)$ se partage en $6$ classes (orbites sous $\mathfrak{S}_{I_0}$),}
\note{ces deux lignes sont biffées d'un trait chacune}

\page{5}
\note{numéroté \emph{4)} en haut à gauche}

\begin{tikzcd}[column sep=small]
\mathbb{D}_{I_0} \subset \mathfrak{S}_{I_0} \ [\subset \mathfrak{S}_I] & \\
\mathrm{C}^{!}_{\mathrm{aff}\,\mathrm{st}}(I_0) \subset \mathrm{C}_{\mathrm{aff}\,\mathrm{st}}(I_0) \arrow[d, "\alpha"] \arrow[dr, "\alpha^{!}"'] & \\
\mathrm{C}_{\mathrm{aff}}(n) & \mathrm{C}_{\mathrm{aff}}(n) \simeq \mathrm{C}^{!}_{\mathrm{aff}\,\mathrm{st}}(I_0) / \mathbb{D}_{I_0}
\end{tikzcd}

\note{schéma redressé : sur la page, une seule fois $\mathrm{C}_{\mathrm{aff}}(n)$, au bas d'une flèche verticale $\alpha$ issue de $\mathrm{C}_{\mathrm{aff}\,\mathrm{st}}(I_0)$ et d'une flèche oblique $\alpha^{!}$ partant de $\mathrm{C}^{!}_{\mathrm{aff}\,\mathrm{st}}(I_0)$ ; nous l'avons dédoublé pour ne pas faire se croiser les flèches. En marge, à droite, sous $\mathbb{D}_{I_0}$ : « (\ill{} automorphismes de la str.\ polygonale $I_0$) », qui glose $\mathbb{D}_{I_0}$ (son D ajouré, rendu $\mathbb{D}$, le groupe diédral). L'indice $\mathrm{st}$ de $\mathrm{C}^{!}_{\mathrm{aff}\,\mathrm{st}}(I_0)/\mathbb{D}_{I_0}$ est \uncertain{lu}}

Soit $\xi^{q} \in \mathrm{C}_{\mathrm{aff}}(n)$, alors $(\alpha^{!})^{-1}(\xi^{q})$ (\uncertain{resp.} $\alpha^{-1}(\xi^{q})$) est une orbite sous $\mathbb{D}_{I_0}$ (\uncertain{resp.} sous $\mathfrak{S}_{I_0}$), et on a sans doute
\[
\alpha^{-1}(\xi^{q}) \xleftarrow{\ \sim\ } \mathfrak{S}_{I_0} \wedge^{\mathbb{D}_{I_0}} (\alpha^{!})^{-1}(\xi^{q})
\]
\struck{id.} d'où
\[
\operatorname{Card} \alpha^{-1}(\xi^{q}) = \Bigl( \frac{n!}{2n} \Bigr) \cdot \operatorname{Card} (\alpha^{!})^{-1}(\xi^{q}),
\qquad
\underbrace{\frac{n!}{2n}}_{} = \frac{(n-1)!}{2}
\]
$= \mathrm{nb}$ de structures polygonales sur $|I_0|$
\note{l'exposant de $\xi$, tout au long des pages 5 à 7, est un petit signe en forme de « q » ou de « 4 » ; nous le rendons $\xi^{q}$ sans être sûr de la lecture. Sur la flèche $\xleftarrow{\sim}$ une première écriture de la formule est raturée et illisible. L'accolade sous $n!/2n$ renvoie à $(n-1)!/2$, puis à « $=$ nb de structures polygonales sur $|I_0|$ »}

\page{6}
\note{numéroté \emph{5)} en haut à gauche}
\[
\operatorname{card} (\alpha^{!})^{-1}(\xi^{q}) = \frac{\operatorname{card} \mathbb{D}_{I_0}}{\operatorname{card} \bigl( G_{\xi} = \operatorname{Aut}(X, \Sigma, \Delta) \bigr)} = \frac{2n}{N(\xi)}
\]
\note{au dénominateur de droite, un premier signe est raturé devant $N(\xi)$}

\uncertain{Prenant} \uncertain{son} composé

\begin{tikzcd}
\widetilde{\mathrm{C}}_{\mathrm{st}}(S) \arrow[d, "\deg 2"] \arrow[dd, bend left=60, "\widetilde{\alpha}"] \\
\mathrm{C}^{!}_{\mathrm{aff}\,\mathrm{st}}(I_0) \arrow[d, "\alpha^{!}"] \\
\mathrm{C}_{\mathrm{aff}}(n) \simeq \widetilde{\mathrm{C}}_{\mathrm{st}}(S)/\mathbb{D}_S \simeq \mathrm{C}^{!}_{\mathrm{aff}\,\mathrm{st}}(I_0)/\mathbb{D}_{I_0}
\end{tikzcd}

on trouve
\[
\operatorname{Card} \widetilde{\alpha}^{-1}(\xi^{q}) = \frac{4n}{N(\xi)}
\]
où
\[
N(\xi) = \operatorname{card} \operatorname{Aut}(X, \Sigma, \Delta)
\]
où
\[
\xi^{q} = \mathrm{cl}_{\mathrm{aff}}(X, \Sigma, \Delta)
\]

\emph{Ex} 1) $n = \operatorname{card} I_0 = 3$, alors \struck{c} $\operatorname{card} \mathrm{C}_{\mathrm{aff}}(3) = 1$, et $\operatorname{Aut}(X, \Sigma, \Delta) \simeq \mathfrak{S}_3$, de card $6$, donc $\widetilde{\mathrm{C}}_{\mathrm{st}}(S)$ consiste en une seule orbite sous $\mathbb{D}_{I_0}$, d'ordre $\frac{12}{6} = 2$

\page{7}
\note{numéroté \emph{6)} en haut à gauche ; c'est le dernier numéro de la suite commencée page 2}
2) $n = 4$, $\operatorname{card} \mathrm{C}_{\mathrm{aff}}(n) = 1$ i.e.\ $\widetilde{\mathrm{C}}_{\mathrm{st}}(S)$ une seule orbite sous $\mathbb{D}_S$. \struck{\ill{}} On a $\operatorname{Aut}(X, \Sigma, \Delta) \simeq \mathbf{Z}/2\mathbf{Z}$, donc $N(\xi^{q}) \simeq 2$, et
\[
\operatorname{card} \widetilde{\mathrm{C}}_{\mathrm{st}}(S) = \frac{4n}{2} = 8
\]
3) $n = 5$, $\operatorname{card} \mathrm{C}_{\mathrm{aff}}(n) = 6 \ldots$

\emph{Résumé}
\[
\widetilde{\mathrm{C}}_{\mathrm{st}}(S) \overset{\mathrm{def}}{=}
\]
classes d'iso.\ de disques avec syst.\ \add{stables} de cordes, épinglés par le polygone comb.\ $S$ d'ordre $2n$.
\note{un double trait horizontal sépare les exemples du « Résumé » ; « stables » est ajouté au-dessus de la ligne}

Dépend fonctoriellement du pol.\ comb.\ $S$ (pour iso comb.), donc $\mathbb{D}_S$ opère dessus. Si $\widetilde{\xi} \in \widetilde{\mathrm{C}}_{\mathrm{st}}(S)$, alors ?
\[
\bigl( \mathbb{D}_S \bigr)_{\widetilde{\xi}} \simeq \operatorname{Aut}_{\mathrm{ist}}(\mathbf{D}, \widetilde{\Sigma}) \simeq \operatorname{Aut}(X, \Sigma, \Delta)
\]
\note{devant $(\mathbb{D}_S)_{\widetilde{\xi}}$, le mot « Aut » est biffé. Sous ce terme : « $=$ stabilisateur de $\widetilde{\xi}$ dans $\mathbb{D}_S$ » ; sous $\operatorname{Aut}_{\mathrm{ist}}$ : « (\ill{} des cat.\ isotopiques) ». L'indice « ist » est \uncertain{lu} ainsi, pour « isotopique »}
(où $(\mathbf{D}, \widetilde{\Sigma}) \longleftrightarrow (X, \Sigma, \Delta)$, et leur classe est $\widetilde{\xi}$). Enfin \ldots
\[
\widetilde{\mathrm{C}}_{\mathrm{st}}(S) / \mathbb{D}_S \simeq \mathrm{C}_{\mathrm{aff}\,\mathrm{st}}(n) \overset{\mathrm{def}}{=}
\]
classes d'iso.\ des syst.\ stables $(X, \Sigma, \Delta)$
\[
\widetilde{\mathrm{C}}_{\mathrm{st}}(S) / \lbrace 1, a_S \rbrace \simeq \mathrm{C}^{!}_{\mathrm{aff}\,\mathrm{st}}(I_0) \quad \bigl( \subset \mathrm{C}_{\mathrm{aff}\,\mathrm{st}}(I_0) \bigr)
\]
où $I_0 = S / \lbrace 1, a_S \rbrace$
\note{l'inclusion $\subset \mathrm{C}_{\mathrm{aff}\,\mathrm{st}}(I_0)$ est écrite au-dessus, dans une bulle rattachée à $\mathrm{C}^{!}_{\mathrm{aff}\,\mathrm{st}}(I_0)$ ; sous cette dernière, un indice griffonné est illisible}

\page{9}
\note{page entièrement occupée par une figure, sans texte suivi : un grand arrangement de pseudo-droites dans un disque (cercle au crayon). Les droites portent, à leurs extrémités sur le bord, les étiquettes appariées $\Delta_A, \Delta_{A'}$, $\Delta_B, \Delta_{B'}$, $\Delta_C, \Delta_{C'}$ (des droites) et $D_A$, $D_B$, $D_C$ (trois courbes qui s'enroulent autour du centre) ; les points d'intersection sont nommés $A, B, C$, $A', B', C'$ au centre, puis $L, M, N$, $L', M', N'$, $P, Q, R$, $P', Q', R'$ plus loin. Chaque région porte le nombre de ses côtés, $3$, $4$, $5$ ou $3'$, $4'$, $5'$ pour les régions symétriques ; à l'infini, les régions entre $\Delta_A$ et $\Delta_{A'}$ (et de même pour $B$, $C$) sont marquées $4'$, et un arc épaissi du bord ferme chacune d'elles. Au-dessous, deux petits hexagones de sommets $A, B', C, A', B, C'$ dans cet ordre, le premier couvert de diagonales et de ratures, le second avec ses trois grandes diagonales}

Décompte des régions, en bas à droite :

triangles : $1 + 3 + 3 = 7$

carrés : $3 + 3 + 3 + 3 + 3 + 3 + 3 + 3 = 24$

pentagones : $3 + 3 = 6$
\note{les mots « triangles », « carrés », « pentagones » sont de nous : sur la page, chaque ligne est précédée d'un petit dessin, triangle, carré, pentagone ; à gauche, des chiffres « 3 », « 3 », « 2 » griffonnés sont en partie raturés. Deux accolades sous la ligne des carrés groupent certains « $3$ »}

\page{10}
\note{changement de papier et de sujet : feuille blanche, écrite au stylo, sans numéro de sa main ; il s'agit de la factorisation de Stein, sans lien apparent avec les pseudo-droites des pages 2 à 9}
Donc si $\widetilde{Z}' \simeq X'$ i.e.\ $Z' \to X'$ est \emph{steinien}
\[
\exists !\ X' \longrightarrow \widetilde{Z}
\]
\note{sous la flèche, en petit et souligné : « $X$-\ill{} » ; le $X'$ source récrit un premier signe biffé}
tel que

\begin{tikzcd}
Z \arrow[d] & Z' \arrow[l] \arrow[d] \\
\widetilde{Z} & X' \arrow[l]
\end{tikzcd}

\note{le carré porte « comm. » en son milieu ; au-dessous, un prolongement vers $X$ (une flèche verticale et une oblique issue de $X'$) est raturé}

En d'autres termes, dans la catégorie des flèches \emph{steiniennes} $Z' \to X'$ \struck{de $\mathrm{Top}/X$} au-dessus de $Z \to X$, il y a un objet initial, savoir $(Z \to \widetilde{Z})$ au-dessus de $(Z \to X)$.

Compatibilités avec chgt de base

Quand peut-on dire que

\begin{tikzcd}
Z \arrow[d] & Z' \arrow[l] \arrow[d] \\
X & X' \arrow[l]
\end{tikzcd}

\uncertain{étant} cartésien, que le carré \ill{} (\uncertain{ou} les deux) de

\begin{tikzcd}
Z \arrow[d] & Z' \arrow[l] \arrow[d] \\
\widetilde{Z} \arrow[d] & \widetilde{Z}' \arrow[l] \arrow[d] \\
X & X' \arrow[l]
\end{tikzcd}

le sont ?

a) $Z \to X$ \uncertain{compact} ?

b) $X' \to X$ loc.\ $\mathcal{O}$-acyclique \ldots
\note{le « $\mathcal{O}$ » est un rond appuyé, d'une autre encre, \uncertain{lu} ainsi}

\page{11}
\note{page de dessins au crayon, avec quelques calculs : une sphère parcourue de trois grands cercles (points d'intersection marqués) ; une pyramide et une bipyramide ; un réseau oblique de points marqués alternativement « $+$ » et « $\bullet$ » ; un octaèdre repassé à l'encre bleue, inscrit dans un cube au crayon, avec à côté un signe \uncertain{$\mathfrak{S}$} ; un cube quadrillé ; en bas à droite un polyèdre à faces carrées et hexagonales (un octaèdre tronqué ?) inscrit dans un cube ; en haut à droite un arbre ramifié à trois niveaux}
\[
48 = 2 \cdot 4! \qquad \frac{5!}{2 \cdot 4!} = \frac{5}{2}
\]
\[
3 \times 2 \times 3 \qquad 12 = \frac{8 \times 3 \times 2}{4}
\]
\note{le « $48$ » récrit un premier chiffre ; « $3 \times 2 \times 3$ » suit, en oblique, un griffonnage raturé sous l'arbre}

\page{12}
\note{un seul dessin au crayon, en haut de la page : un arrangement de six pseudo-droites, dont les extrémités portent des indices $1, 2, 3, 4, 5$ précédés du même signe que les régions. Au centre, un pentagone marqué $5$, entouré de cinq triangles marqués $3$ ; autour, six régions marquées $4_1, \ldots, 4_6$ (le chiffre « $4$ » a la même forme que l'exposant de $\xi$ aux pages 5 à 7). Au-dessus, à gauche, deux traits « ${}^{\prime\prime}$ » isolés}

\page{13}
\note{feuille à l'italienne, au stylo et au crayon. En haut, deux petites figures : des carrés recoupés par des droites obliques, certaines en pointillé (à gauche, deux coins marqués « ${}^{\prime\prime}$ »). En bas, une construction d'un arrangement de droites : deux horizontales, deux verticales et des obliques, avec les points $P$, $L$, $B$, $C'$, $N$ sur la première horizontale, $N'$, $C$, $B'$, $L'$, $P'$ sur la seconde, $A$ et $A'$, $Q$ et $Q'$, $M$ et $M'$, $R$ et $R'$ ; les longueurs $a$, $b$ (en haut) et $a'$, $b'$ (en bas) sont marquées sur les segments. À droite de la figure :}
\[
a + a' \leqslant 2 , \qquad b + b' \geqslant 2 ,
\qquad 2 - a' < b \quad \text{i.e.} \quad a' + b > 2
\]
(d'\uncertain{où} $b > a$)
\note{dans les deux premières inégalités le signe a été repassé à l'encre, sans qu'on puisse dire s'il s'agit de l'inégalité large ou stricte}

\page{14}
\note{feuille à l'italienne, dessins au stylo sans texte : un arrangement de cinq droites $D_0$ (horizontale, avec une variante $D_0'$ \uncertain{lue}, repassée), $D_1$, $D_2$, $D_3$, $D_4$ formant une étoile, deux triangles inférieurs hachurés et le demi-plan sous $D_0$ parcouru de traits parallèles ; à droite, une courbe coupe l'étoile (pseudo-droite), et quatre petits croquis de droites concourantes ou sécantes, le côté de chaque demi-plan marqué par des hachures}

\section{Système d'hyperplans (projectifs) réels}
\note{titre écrit de sa main en tête de la page 16, souligné ; la page n'est pas numérotée par lui, et reprend sur papier blanc au stylo noir}

\page{16}
\marginal{Théorie analogue dans le cas affine, \ill{} \ill{} mais qui ne \uncertain{passent} pas \ill{} par \ill{} \ill{} \ill{} des $H_i$ \ill{}}
Soit $V$ esp.\ vectoriel réel, $H$ un hyperplan de $V$, on note $D_H$ \add{$\subset \mathfrak{P}(V)$} l'ens des trois éléments \struck{de} $\lbrace H, V'_H, V''_H \rbrace$, où $V'_H$ et $V''_H$ sont les deux demi-espaces ouverts de $V$ déterminés par $H$ ; \struck{\uncertain{quitte}} donc $D_H$ est une partition de $V$ par \uncertain{des} \add{ens} convexes.
\note{la note marginale est écrite en oblique dans le coin supérieur gauche ; « $\subset \mathfrak{P}(V)$ » est ajouté au-dessus de $D_H$ et rattaché par un trait ; devant $D_H$ un sigle est raturé}

Comme c'est un ensemble à trois éléments, avec un élément distingué $\lbrace H \rbrace$, on peut regarder $D_H$ comme un \uncertain{espace} isomorphe à $\mathbf{F}_3 = \lbrace -1, 0, 1 \rbrace$. On a une application canonique \struck{\ill{} $V \setminus$}
\[
\varphi_H : \mathbf{S} \xrightarrow{\ \varphi_H\ } D_H , \qquad \mathbf{S} = V^{*} / \mathbf{R}^{*}_{+}
\]
\note{son $\mathbf{S}$ est un S doublé ; au-dessous, « $= V^{*}/\mathbf{R}^{*+}$ », écrit $\mathbf{R}^{*+}$ sur la page. Ici $V^{*}$ désigne \uncertain{sans doute} $V \setminus \lbrace 0 \rbrace$, ce qui fait de $\mathbf{S}$ la sphère des directions}

Soit $(H_i)_{i \in I}$ une famille finie d'hyperplans de $V$, \emph{distincts}, et considérons l'application
\marginal{« distincts » \ill{} \ill{} \ill{} $\varphi$ \ill{} \ill{} \ill{}}
\[
\varphi : \mathbf{S} \longrightarrow \prod D_{H_i} = \mathbf{E}
\]
définie par les $\varphi_{H_i}$. Les fibres non vides de $\varphi$ forment une partition de $S$ par des \struck{\ill{}} \add{parties} « convexes » (i.e.\ l'image inverse dans $V^{*}$ \struck{\ill{}} est un cône convexe) \emph{dans} les « alvéoles ouvertes » :
\marginal{« \uncertain{convexes} » \ill{} \ill{} \ill{} \ill{} \ill{} \ill{} question \ill{} \ill{} \ill{} \ill{} si $\bigcap H_i = \lbrace 0 \rbrace$ \ill{} \ill{} si $\bigcap H_i = \lbrace 0 \rbrace$}
\note{la note marginale d'en bas est écrite en oblique sur cinq ou six lignes serrées, et un trait la rattache au mot « \uncertain{convexes} » ; « si $\bigcap H_i = \lbrace 0 \rbrace$ » s'y lit deux fois. Le passage « (i.e.\ \ldots\ $\lbrace 0 \rbrace$, est un cône convexe) » porte au-dessus de « cône » un « un » ajouté ; le $\lbrace 0 \rbrace$ commence une ligne et le reste de la phrase est \uncertain{reconstitué}}
\[
\text{Soit } \Sigma = \varphi(\mathbf{S}) \subset \mathbf{E}
\]
l'ens.\ des indices de cette partition. C'est une partie de $\mathbf{E}$ stable par l'homothétie $-1$ ($-\mathrm{id}_{\mathbf{E}}$), qui correspond par passage \struck{\ill{}} d'une fibre
\[
Z^{*}_{-\alpha} = - Z^{*}_{\alpha} \qquad \text{pour } \alpha \in \Sigma
\]
On se propose de déterminer $\overline{Z}^{*}_{\alpha}$ pour $i \in I$.
\note{devant $\Sigma = \varphi(\mathbf{S})$ un premier symbole est raturé ; dans « pour $\alpha \in \Sigma$ » le $\Sigma$ est souligné. Dans la dernière phrase, « pour $i \in I$ » est lu tel qu'il est écrit, où l'on attendrait $\alpha \in \Sigma$ ; l'adhérence $\overline{Z}^{*}_{\alpha}$ est calculée page 17}

\page{17}
\marginal{\ill{} \ill{} \ill{} \ill{} $H_i$ \ill{} \ill{} \ill{}}
\note{la note marginale du coin supérieur gauche, en oblique, est barrée de plusieurs grands traits en boucle ; elle est illisible}
\ill{} trouve que \uncertain{via} l'\uncertain{intersection} des $\frac{1}{2}$ espaces \emph{fermés}
\[
\overline{V'_{H_i}} = V'_{H_i} \cup H_i , \qquad \overline{V''_{H_i}} = V''_{H_i} \cup H_i ,
\]
déterminés par les $H_i$ ($i \in I$), \add{qui contiennent $Z_i$}, en d'autres termes
\[
\overline{Z}^{*}_{\alpha} = \bigcup_{\substack{\beta \in \Sigma \\ \text{tel que } i \in I \\ \text{implique } k_i(\beta) \leqslant k_i(\alpha)}} Z^{*}_{\beta}
= \bigcup_{\substack{\beta \in \Sigma \\ \beta \preceq \alpha}} Z^{*}_{\beta}
\]
\note{« $\frac{1}{2}$ espaces » : il écrit « ½ » pour « demi- », cf.\ l'aide-lecture. « qui contiennent $Z_i$ » est ajouté au-dessus de la ligne, dans une bulle ; l'indice de $Z$ y est \uncertain{lu} $i$. Dans l'inégalité $k_i(\beta) \leqslant k_i(\alpha)$ le signe récrit un premier signe raturé}
où les \struck{\ill{}}
\[
k_i : \mathbf{E} \longrightarrow \mathbb{D}_{H_i} \qquad (i \in I)
\]
sont les fonctions coordonnées, et où on introduit dans $\mathbb{D}_{H_i}$ la relation d'ordre $\leqslant$ \uncertain{définie} par

\begin{tikzcd}[column sep=small]
\alpha'_i & & \alpha''_i \\
& 0 \arrow[ul, no head] \arrow[ur, no head] &
\end{tikzcd}

$\alpha', \alpha''$ les deux éléments $\neq 0$ de $\mathbb{D}_{H_i}$,
\note{ici, et jusqu'au bas de la page, il écrit $\mathbb{D}_{H_i}$ (un D ajouré) pour l'ensemble à trois éléments noté $D_{H_i}$ à la page 16 ; le petit schéma est un « V » dont la pointe est marquée $0$}
i.e.\ si $\xi, \eta \in \mathbb{D}_{H_i}$, on a $\xi \leqslant \eta$ ssi $\xi = 0$ ou $\xi = \eta$, et on met sur $\mathbf{E} = \prod_{i \in I} \mathbb{D}_i$ l'ordre produit. \struck{\ill{} est ordonné}, \uncertain{puis} sur $\Sigma$ l'ordre induit, d'où une « top.\ inférieure » sur $\Sigma$, dont les \emph{fermés} sont les parties $A$ de $\Sigma$ telles que $\alpha \in A$, $\beta \leqslant \alpha \Rightarrow \beta \in A$. On a donc
\[
\beta \leqslant \alpha \iff \beta \in \overline{\lbrace \alpha \rbrace} , \quad \text{et} \quad
\overline{Z}^{*}_{\alpha} = Z^{*}_{\overline{\lbrace \alpha \rbrace}}
\]
où on pose, pour $A \subset \Sigma$,
\[
Z^{*}_{A} = \bigcup_{\alpha \in A} Z^{*}_{\alpha} .
\]

\page{18}
Plus généralement, on a
\[
\overline{Z^{*}_{A}} = Z^{*}_{\overline{A}}
\]
\[
\Bigl( \bigcup_{\alpha \in A} \overline{Z}_{\alpha} = \bigcup_{\alpha \in A} Z_{\overline{\lbrace \alpha \rbrace}} = Z_{\bigcup_{\alpha \in A} \overline{\lbrace \alpha \rbrace}} , \quad \bigcup_{\alpha \in A} \overline{\lbrace \alpha \rbrace} = \overline{A} \Bigr)
\]
\note{la première union est reliée à $\overline{Z^{*}_{A}}$ par un « $=$ » vertical ; sous la dernière, une accolade donne $\overline{A}$}

L'ens.\ \struck{$\Sigma$} ordonné $\Sigma$ est \struck{donc} \add{donc} (\uncertain{fini}) l'ens.\ d'indices d'une stratification de la sphère $S$, satisfaisant l'axiome de \emph{frontière} \add{et de densité} (\struck{\ill{} \ill{}} $\Sigma_\alpha$ \struck{\ill{} \ill{}}) (les $Z^{*}_{\alpha} \neq \emptyset$, $\overline{Z^{*}_{\alpha}} = Z_{\overline{\alpha}}$ \struck{\ill{}}), et à strates convexes.
\note{« donc » est récrit au-dessus de « donc » biffé ; « et de densité » est ajouté en marge et renvoyé dans la ligne ; la parenthèse biffée contient un $\Sigma_\alpha$ repris au-dessus, lui aussi raturé. Le dernier terme de l'égalité est \uncertain{lu} $Z_{\overline{\alpha}}$}

\struck{$\Sigma$ en général n'est pas \struck{un \ill{}} \emph{Inf}-\uncertain{réticulé} : \uncertain{deux} éléments de $\Sigma$ qui ont minorants n'ont pas nécessairement de Inf -- ex.\ $V$ de dim $3$, \ill{} \ill{} plans passant par une même droite ($n \geqslant 3$ p.\ ex.\ $n = 3$). Mais sans \uncertain{doute}, si l'intersection des $H_i$ se réduit à $\lbrace 0 \rbrace$, \struck{\ill{}} l'intersection de deux fermés \struck{\ill{}} fermés sera une partie fermée.}
\note{ce paragraphe est barré par cinq longs traits obliques, et sa dernière ligne d'un trait horizontal. En regard, dans la marge gauche, un petit croquis : une sphère (ou un disque) où deux arcs se coupent en deux points antipodaux, marqué d'un « ? »}

$\Sigma$ est une partie ouverte de $\mathbf{E}$, i.e.\ $\alpha \in \Sigma$, $\beta \in \mathbf{E}$, $\beta \geqslant \alpha$ implique $\beta \in \Sigma$. Donc dans $\Sigma$ comme dans $\mathbf{E}$ (\uncertain{et} partie \add{de} \uncertain{même} \uncertain{genre} \ill{} une \uncertain{façon} canonique \struck{dans} $\mathbf{F}_3$) toute partie \uncertain{minorée} de $\Sigma$ a un Inf.
\note{la fin de la page est d'une lecture très incertaine ; « Inf » est écrit « I--f », avec une longue queue}

\page{20}
\note{feuille à l'italienne. En haut, quatre dessins d'arrangements de droites, séparés par un double trait vertical après le premier : (1) quatre droites en position générale, marqué « $4$ » et « ($\mathfrak{S}_4$) » ; (2) cinq droites, « ($D_5$) $5$ pos.\ générale » ; (3) cinq droites, « $5$ position particulière » ; (4) six droites, les trois côtés d'un triangle et ses trois médianes concourantes, points d'intersection marqués, « $6$ en position \uncertain{particulière} ». Entre (2) et (3), un signe raturé}
Tout arr.\ \uncertain{que} $1$-dim.\ de ps.\ droites \struck{\ill{}} pouvant \uncertain{être} de ps.\ droites \emph{bigones}, \struck{contient un des $4$ configurations} \add{s'il ne se réduit à $4$ ps.\ droites} (\uncertain{auquel} cas c'est le cas (1)) contient une des \add{types de} configurations (2), (3), (4) qui sont \struck{spéciales} \add{2-dim.} « non spéciales », i.e.\ quand elles \uncertain{sont} réalisées \add{\uncertain{comme}} des arrangements (de ps.\ droites, dans l'intersection de deux faces fermées est \struck{\ill{}} vide, réduite à un sommet, ou est une arête fermée.
\note{deux corrections interlinéaires : « s'il ne se réduit à $4$ ps.\ droites » est écrit au-dessus de « contient un des $4$ configurations » biffé ; « types de » est ajouté au-dessus de « configurations ». « 2-dim. » est écrit au-dessous de « spéciales » biffé, et un grand rond, relié par un trait, renvoie à « arrangement ». La parenthèse ouverte avant « de ps.\ droites » n'est pas refermée}

\emph{NB} Si elle ne contient (2) ni (3), \struck{c'est} elle est du type (4), ou du type (5)
\note{la configuration (5), dessinée à droite, est un arrangement de sept droites, dont trois en tirets, avec six points d'intersection marqués}
des $7$ ps.\ droites formé d'un \uncertain{quadrilatère} et de ses \uncertain{trois} diagonales, qui \uncertain{mérite} d'être étudié.
\note{la phrase court sur deux colonnes, de part et d'autre du dessin ; l'ordre de lecture retenu est le nôtre}

\struck{Si elle ne contient}
\note{la page s'arrête sur ces mots biffés ; la suite est dans le lot suivant, s'il y en a une}

\end{document}
